\documentclass[11pt]{scrartcl}

\usepackage[top=1.5cm]{geometry}
\usepackage{float}
\usepackage{listings}
\usepackage{xcolor}
\usepackage{hyperref}

\setlength{\parindent}{0em}
\setlength{\parskip}{0.5em}

\newcommand{\youranswerhere}{[Your answer goes here \ldots]}
\renewcommand{\thesubsection}{\arabic{subsection}}

\lstdefinestyle{dbtsql}{
  language=SQL,
  basicstyle=\small\ttfamily,
  keywordstyle=\color{magenta!75!black},
  stringstyle=\color{green!50!black},
  showspaces=false,
  showstringspaces=false,
  commentstyle=\color{gray}}

\title{
  \textbf{\large Assignment 3} \\
  Index Tuning -- Selection \\
  {\large Database Tuning}}

\author{
  Group B15 \\
  \large El Lbidi Omar, 12046411 \\
  \large Lastname2 Firstname2, StudentID2 \\
  \large Lastname3 Firstname3, StudentID3
}

\begin{document}

\maketitle

\paragraph{Notes}

\begin{itemize}
  \item Do not forget to run \lstinline[style=dbtsql]{ANALYZE tablename} after creating or changing a table.
  \item Use \lstinline[style=dbtsql]{EXPLAIN ANALYZE} for the query plans that you display in the report.
\end{itemize}

\subsection*{Experimental Setup}

How do you send the queries to the database? How do you measure the execution time for a sequence of queries?

We sent queries to the PostgreSQL database using a Python script (\texttt{test\_baseline.py}) connecting via a database driver. To avoid network latencies skewing the results, the client script was executed directly on the local database server.

To measure execution time, we utilized a \texttt{while} loop constrained by the system clock (\texttt{time.time() - start\_time < 65}). This continuously executed the target query for a fixed duration of 65 seconds. A counter tracked the total number of queries completed. The average runtime per query was calculated by dividing the total duration by the total number of queries. Throughput (queries per second) is calculated as the inverse of the average runtime.

\textbf{Important Note on Caching:} To prevent database caching from artificially inflating our throughput, we used Python's \texttt{random.choice()} function to select a different condition from a pre-defined list of 10 values for every single query iteration. If we had repeated the exact same condition, PostgreSQL would have served the result from its memory cache, and we would be measuring RAM speed rather than true index traversal performance.

\subsection*{Clustering B$^+$ Tree Index}

\paragraph{Point Query}

Repeat the following query multiple times with different conditions for \texttt{pubID}.

\begin{lstlisting}[style=dbtsql]
SELECT * FROM Publ WHERE pubID = ...
\end{lstlisting}

Which conditions did you use?

We randomly cycled through 10 distinct pubIDs: \\
\texttt{'journals/automatica/DehghaniLA', 'conf/crypto/BressonSS02', 'conf/isbi/ArdekaniS04', 'conf/dac/Lanzo92', 'conf/ipps/Sobe04', 'conf/conpar/LeiteB92', 'journals/tomccap/GhineaDGC07', 'conf/cp/SimonisO08', 'journals/jal/Xu00', 'conf/sigmod/LiuC07'}

Show the runtime results and compute the throughput.

Average Runtime: 0.000045 seconds \\
Throughput: $\approx$ 22,222 queries/second

Query plan (for one of the queries):

{\small
\parskip0pt\begin{verbatim}
Index Scan using pubid_bplus_idx on publ (cost=0.43..8.45 rows=1 width=112)
  Index Cond: ((pubid)::text = 'journals/automatica/DehghaniLA'::text)
\end{verbatim}}

\paragraph{Multipoint Query vs. Multipoint Query IN-Predicate -- Low Selectivity}

Repeat the following query multiple times with different conditions for \texttt{booktitle}.

\begin{lstlisting}[style=dbtsql]
SELECT * FROM Publ WHERE booktitle = ...
\end{lstlisting}

\begin{lstlisting}[style=dbtsql]
SELECT * FROM Publ WHERE pubID IN (...)
\end{lstlisting}


Which conditions did you use?

For \texttt{booktitle}, we randomly cycled through 10 distinct titles: \\
\texttt{'SBRN', 'WRAC', 'IJCSLP Workshop...', 'ICVS', 'RelMiCS', 'Security and Management', 'RSKD', 'Computer Science Today', 'Architecture and Protocols...', 'Computational Discovery...'}

For the IN-predicate, we randomly sampled 3 unique \texttt{pubID}s per query from our pool of 10 distinct pubIDs.

Show the runtime results and compute the throughput.

Low Selectivity (\texttt{booktitle}): \\
Average Runtime: 0.000126 seconds \\
Throughput: $\approx 7,936$ queries/second

\vspace{0.5em}
IN-Predicate (\texttt{pubID}): \\
Average Runtime: 0.000068 seconds \\
Throughput: $\approx 14,705$ queries/second

Query plan (for one of the queries):

{\small
\parskip0pt\begin{verbatim}
Index Scan using title_bplus_idx on publ (cost=0.43..14.54 rows=178 width=112)
  Index Cond: ((booktitle)::text = 'SBRN'::text)
\end{verbatim}}

\paragraph{Multipoint Query -- High Selectivity}

Repeat the following query multiple times with different conditions for \texttt{year}.

\begin{lstlisting}[style=dbtsql]
SELECT * FROM Publ WHERE year = ...
\end{lstlisting}

Which conditions did you use?

We randomly cycled through 10 distinct years: \\
\texttt{'1974', '1952', '1994', '1959', '2005', '1976', '1950', '1991', '1968', '1987'}

Show the runtime results and compute the throughput.

Average Runtime: 0.010828 seconds \\
Throughput: $\approx 92$ queries/second

Query plan (for one of the queries):

{\small
\parskip0pt\begin{verbatim}
Index Scan using year_bplus_idx on publ (cost=0.43..87.11 rows=2096 width=112)
  Index Cond: ((year)::text = '1974'::text)
\end{verbatim}}

\subsection*{Non-Clustering B$^+$ Tree Index}

\emph{Note:} Make sure the data is not physically ordered by the indexed attributes due to the clustering index that you created before. \\
\emph{To ensure the data was not physically ordered by the attribute being tested, we physically scattered the data prior to the test by clustering the table on an independent attribute (e.g., clustering by \texttt{booktitle} before testing the index on \texttt{year}).}

\paragraph{Point Query}

Repeat the following query multiple times with different conditions for \texttt{pubID}.

\begin{lstlisting}[style=dbtsql]
SELECT * FROM Publ WHERE pubID = ...
\end{lstlisting}

Which conditions did you use?

We randomly cycled through 10 distinct pubIDs: \\
\texttt{'journals/automatica/DehghaniLA', 'conf/crypto/BressonSS02', 'conf/isbi/ArdekaniS04', 'conf/dac/Lanzo92', 'conf/ipps/Sobe04', 'conf/conpar/LeiteB92', 'journals/tomccap/GhineaDGC07', 'conf/cp/SimonisO08', 'journals/jal/Xu00', 'conf/sigmod/LiuC07'}

Show the runtime results and compute the throughput.

Average Runtime: 0.000046 seconds \\
Throughput: $\approx 21,739$ queries/second

Query plan (for one of the queries):

{\small
\parskip0pt\begin{verbatim}
Index Scan using pubid_bplus_idx on publ (cost=0.43..8.45 rows=1 width=112)
  Index Cond: ((pubid)::text = 'journals/automatica/DehghaniLA'::text)
\end{verbatim}}

\paragraph{Multipoint Query vs. Multipoint Query IN-Predicate -- Low Selectivity}

Repeat the following query multiple times with different conditions for \texttt{booktitle}.

\begin{lstlisting}[style=dbtsql]
SELECT * FROM Publ WHERE booktitle = ...
\end{lstlisting}

\begin{lstlisting}[style=dbtsql]
SELECT * FROM Publ WHERE pubID IN (...)
\end{lstlisting}


Which conditions did you use?

For \texttt{booktitle}, we randomly cycled through 10 distinct titles: \\
\texttt{'SBRN', 'WRAC', 'IJCSLP Workshop...', 'ICVS', 'RelMiCS', 'Security and Management', 'RSKD', 'Computer Science Today', 'Architecture and Protocols...', 'Computational Discovery...'}

For the IN-predicate, we randomly sampled 3 unique \texttt{pubID}s per query from our pool of 10 distinct pubIDs.

Show the runtime results and compute the throughput.

Low Selectivity (\texttt{booktitle}): \\
Average Runtime: 0.000126 seconds \\
Throughput: $\approx 7,936$ queries/second

\vspace{0.5em}
IN-Predicate (\texttt{pubID}): \\
Average Runtime: 0.000069 seconds \\
Throughput: $\approx 14,492$ queries/second

Query plan (for one of the queries):

{\small
\parskip0pt\begin{verbatim}
Bitmap Heap Scan on publ (cost=5.80..668.84 rows=177 width=112)
  Recheck Cond: ((booktitle)::text = 'SBRN'::text)
  -> Bitmap Index Scan on title_bplus_idx (cost=0.00..5.75 rows=177 width=0)
       Index Cond: ((booktitle)::text = 'SBRN'::text)
\end{verbatim}}

\paragraph{Multipoint Query -- High Selectivity}

Repeat the following query multiple times with different conditions for \texttt{year}.

\begin{lstlisting}[style=dbtsql]
SELECT * FROM Publ WHERE year = ...
\end{lstlisting}

Which conditions did you use?

We randomly cycled through 10 distinct years: \\
\texttt{'1974', '1952', '1994', '1959', '2005', '1976', '1950', '1991', '1968', '1987'}

Show the runtime results and compute the throughput.

Average Runtime: 0.012559 seconds \\
Throughput: $\approx 79$ queries/second

Query plan (for one of the queries):

{\small
\parskip0pt\begin{verbatim}
Bitmap Heap Scan on publ (cost=52.42..5044.40 rows=1603 width=112)
  Recheck Cond: ((year)::text = '1974'::text)
  -> Bitmap Index Scan on year_bplus_idx (cost=0.00..52.02 rows=1603 width=0)
       Index Cond: ((year)::text = '1974'::text)
\end{verbatim}}

\subsection*{Non-Clustering Hash Index}

\emph{Note:} Make sure the data is not physically ordered by the indexed attributes due to the clustering index that you created before. \\
\emph{To ensure the data was not physically ordered by the attribute being tested, we physically scattered the data prior to the test by clustering the table on an independent attribute (e.g., clustering by \texttt{booktitle} before testing the index on \texttt{year}).}

\paragraph{Point Query}

Repeat the following query multiple times with different conditions for \texttt{pubID}.

\begin{lstlisting}[style=dbtsql]
SELECT * FROM Publ WHERE pubID = ...
\end{lstlisting}

Which conditions did you use?

We randomly cycled through 10 distinct pubIDs: \\
\texttt{'journals/automatica/DehghaniLA', 'conf/crypto/BressonSS02', 'conf/isbi/ArdekaniS04', 'conf/dac/Lanzo92', 'conf/ipps/Sobe04', 'conf/conpar/LeiteB92', 'journals/tomccap/GhineaDGC07', 'conf/cp/SimonisO08', 'journals/jal/Xu00', 'conf/sigmod/LiuC07'}

Show the runtime results and compute the throughput.

Average Runtime: 0.000039 seconds \\
Throughput: $\approx 25,641$ queries/second

Query plan (for one of the queries):

{\small
\parskip0pt\begin{verbatim}
Index Scan using pubid_hash_idx on publ (cost=0.00..8.02 rows=1 width=112)
  Index Cond: ((pubid)::text = 'journals/automatica/DehghaniLA'::text)
\end{verbatim}}

\paragraph{Multipoint Query vs. Multipoint Query IN-Predicate -- Low Selectivity}

Repeat the following query multiple times with different conditions for \texttt{booktitle}.

\begin{lstlisting}[style=dbtsql]
SELECT * FROM Publ WHERE booktitle = ...
\end{lstlisting}

\begin{lstlisting}[style=dbtsql]
SELECT * FROM Publ WHERE pubID IN (...)
\end{lstlisting}


Which conditions did you use?

For \texttt{booktitle}, we randomly cycled through 10 distinct titles: \\
\texttt{'SBRN', 'WRAC', 'IJCSLP Workshop...', 'ICVS', 'RelMiCS', 'Security and Management', 'RSKD', 'Computer Science Today', 'Architecture and Protocols...', 'Computational Discovery...'}

For the IN-predicate, we randomly sampled 3 unique \texttt{pubID}s per query from our pool of 10 distinct pubIDs.

Show the runtime results and compute the throughput.

Low Selectivity (\texttt{booktitle}): \\
Average Runtime: 0.000132 seconds \\
Throughput: $\approx 7,575$ queries/second

\vspace{0.5em}
IN-Predicate (\texttt{pubID}): \\
Average Runtime: 0.000046 seconds \\
Throughput: $\approx 21,739$ queries/second

Query plan (for one of the queries):

{\small
\parskip0pt\begin{verbatim}
Bitmap Heap Scan on publ (cost=12.02..23.96 rows=3 width=112)
  Recheck Cond: ((pubid)::text = ANY ('{...}'::text[]))
  -> Bitmap Index Scan on pubid_hash_idx (cost=0.00..12.02 rows=3 width=0)
       Index Cond: ((pubid)::text = ANY ('{...}'::text[]))
\end{verbatim}}

\paragraph{Multipoint Query -- High Selectivity}

Repeat the following query multiple times with different conditions for \texttt{year}.

\begin{lstlisting}[style=dbtsql]
SELECT * FROM Publ WHERE year = ...
\end{lstlisting}

Which conditions did you use?

We randomly cycled through 10 distinct years: \\
\texttt{'1974', '1952', '1994', '1959', '2005', '1976', '1950', '1991', '1968', '1987'}

Show the runtime results and compute the throughput.

Average Runtime: 0.016465 seconds \\
Throughput: $\approx 60$ queries/second

Query plan (for one of the queries):

{\small
\parskip0pt\begin{verbatim}
Bitmap Heap Scan on publ (cost=52.42..5044.40 rows=1603 width=112)
  Recheck Cond: ((year)::text = '1974'::text)
  -> Bitmap Index Scan on year_hash_idx (cost=0.00..52.02 rows=1603 width=0)
       Index Cond: ((year)::text = '1974'::text)
\end{verbatim}}

\subsection*{Table Scan}

\emph{Note:} Make sure the data is not physically ordered by the indexed attributes due to the clustering index that you created before. \\
\emph{To ensure the data was not physically ordered by the attribute being tested, we physically scattered the data prior to the test by clustering the table on an independent attribute (e.g., clustering by \texttt{booktitle} before testing the table scan for \texttt{year}).}

\paragraph{Point Query}

Repeat the following query multiple times with different conditions for \texttt{pubID}.

\begin{lstlisting}[style=dbtsql]
SELECT * FROM Publ WHERE pubID = ...
\end{lstlisting}

\begin{lstlisting}[style=dbtsql]
SELECT * FROM Publ WHERE pubID IN (...)
\end{lstlisting}


Which conditions did you use?

We randomly cycled through 10 distinct pubIDs: \\
\texttt{'journals/automatica/DehghaniLA', 'conf/crypto/BressonSS02', 'conf/isbi/ArdekaniS04', 'conf/dac/Lanzo92', 'conf/ipps/Sobe04', 'conf/conpar/LeiteB92', 'journals/tomccap/GhineaDGC07', 'conf/cp/SimonisO08', 'journals/jal/Xu00', 'conf/sigmod/LiuC07'}

For the IN-predicate, we randomly sampled 3 unique \texttt{pubID}s per query from our pool of 10 distinct pubIDs.

Show the runtime results and compute the throughput.

Point (\texttt{pubID}): \\
Average Runtime: 0.026650 seconds \\
Throughput: $\approx 38$ queries/second

\vspace{0.5em}
IN-Predicate (\texttt{pubID}): \\
Average Runtime: 0.040732 seconds \\
Throughput: $\approx 25$ queries/second

Query plan (for one of the queries):

{\small
\parskip0pt\begin{verbatim}
Gather (cost=1000.00..29785.09 rows=1 width=112)
  Workers Planned: 2
  -> Parallel Seq Scan on publ (cost=0.00..28784.99 rows=1 width=112)
       Filter: ((pubid)::text = 'conf/sigmod/LiuC07'::text)
\end{verbatim}}

\paragraph{Multipoint Query vs. Multipoint Query IN-Predicate -- Low Selectivity}

Repeat the following query multiple times with different conditions for \texttt{booktitle}.

\begin{lstlisting}[style=dbtsql]
SELECT * FROM Publ WHERE booktitle = ...
\end{lstlisting}

\begin{lstlisting}[style=dbtsql]
SELECT * FROM Publ WHERE pubID IN (...)
\end{lstlisting}

Which conditions did you use?

For \texttt{booktitle}, we randomly cycled through 10 distinct titles: \\
\texttt{'SBRN', 'WRAC', 'IJCSLP Workshop...', 'ICVS', 'RelMiCS', 'Security and Management', 'RSKD', 'Computer Science Today', 'Architecture and Protocols...', 'Computational Discovery...'}

Show the runtime results and compute the throughput.

Average Runtime: 0.031166 seconds \\
Throughput: $\approx 32$ queries/second

Query plan (for one of the queries):

{\small
\parskip0pt\begin{verbatim}
Gather (cost=1000.00..29802.79 rows=178 width=112)
  Workers Planned: 2
  -> Parallel Seq Scan on publ (cost=0.00..28784.99 rows=74 width=112)
       Filter: ((booktitle)::text = 'SBRN'::text)
\end{verbatim}}

\paragraph{Multipoint Query -- High Selectivity}

Repeat the following query multiple times with different conditions for \texttt{year}.

\begin{lstlisting}[style=dbtsql]
SELECT * FROM Publ WHERE year = ...
\end{lstlisting}

Which conditions did you use?

We randomly cycled through 10 distinct years: \\
\texttt{'1974', '1952', '1994', '1959', '2005', '1976', '1950', '1991', '1968', '1987'}

Show the runtime results and compute the throughput.

Average Runtime: 0.048356 seconds \\
Throughput: $\approx 21$ queries/second

Query plan (for one of the queries):

{\small
\parskip0pt\begin{verbatim}
Gather (cost=1000.00..29957.59 rows=1726 width=112)
  Workers Planned: 2
  -> Parallel Seq Scan on publ (cost=0.00..28784.99 rows=719 width=112)
       Filter: ((year)::text = '1974'::text)
\end{verbatim}}

\subsection*{Discussion}

Give the throughput of the query types and index types in queries/second.
\begin{table}[H]
  \centering
  \begin{tabular}{c|c|c|c|c}
    & clustering & non-clust.\ B$^+$ tree & non-clust.\ hash & table scan
      \tabularnewline
    \hline
    point (\texttt{pubID}) & 22,222 & 21,739 & 25,641 & 38 \tabularnewline
    \hline
    multipoint (\texttt{booktitle}) & 7,936 & 7,936 & 7,575 & 32
      \tabularnewline
    \hline
    multipoint-IN (\texttt{pubID}) & 14,705 & 14,492 & 21,739 & 25
      \tabularnewline
    \hline
    multipoint (\texttt{year}) & 92 & 79 & 60 & 21
      \tabularnewline
  \end{tabular}
\end{table}

Discuss the runtime results for the different index types and the table scan. Are the results expected? Why (not)?

The results align strongly with database architectural expectations.

For the Point Query, the Hash index significantly outperformed both B$^+$ Trees (25,641 vs $\sim$22,000 q/s). This is expected, as Hash indexes evaluate an exact match in $O(1)$ time, whereas B$^+$ Trees require $O(\log n)$ tree traversal.

For the High Selectivity query (\texttt{year}), the Clustered B$^+$ Tree performed the best (92 q/s) because sequential I/O allowed the engine to read adjacent blocks rapidly. The Non-Clustered Hash index performed the worst here (60 q/s). Hash functions randomly distribute data; thus, finding $\sim$2,000 matching rows required the database to jump to scattered disk addresses, maximizing random I/O penalties.

Interestingly, the difference between Clustered and Non-Clustered B$^+$ Trees for Low Selectivity (\texttt{booktitle}) was negligible. This is due to a combination of PostgreSQL's Bitmap Heap Scan (which buffers and sorts scattered physical addresses in memory to minimize disk head movement) and the high random-access speeds of modern NVMe SSDs. Finally, Table Scans (Sequential Scans) were broadly outclassed across the board as expected, because scanning 3 million rows for a small subset of data is fundamentally inefficient.

\subsection*{Time Spent on this Assignment}

Time in hours per person: \\ 
Omar El Lbidi: 6.5 Hours



\subsection*{References}
\begin{table}[H]
  \centering
  \begin{tabular}{c}
    \hline
    \textbf{Important:} Reference your information sources! \tabularnewline
    Remove this section if you use footnotes to reference your information sources. \tabularnewline
    \hline
  \end{tabular}
\end{table}

\begin{enumerate}
  \item PostgreSQL Global Development Group. \textit{PostgreSQL 14 Documentation: Index Types}. \\
        \href{https://www.postgresql.org/docs/14/indexes-types.html}{https://www.postgresql.org/docs/14/indexes-types.html}
\end{enumerate}

\end{document}